% Осень26!A8; нумерация внутри столбца без пустых строк.
% Источник: лекция 3, §13 (строка 349)
\ticket{06}{Теорема Лагранжа для отображений}
\label{ticket:06}

\begin{theorem}
Пусть $D\subset\mathbb R^m$ открыто, $F:D\to\mathbb R^\ell$ дифференцируемо и $[a,b]\subset D$.
Тогда существует $c\in[a,b]$ такое, что
\[
  |F(b)-F(a)|\le\|F'(c)\|\,|b-a|.
\]
\end{theorem}

\begin{proof}
Положим $h=b-a$, $v=F(b)-F(a)$ и $\psi(t)=\langle F(a+th),v\rangle$.
При $v=0$ утверждение очевидно.
В противном случае по одномерной теореме Лагранжа для некоторого $\theta\in(0,1)$
\[
  |v|^2=\psi(1)-\psi(0)
        =\langle F'(a+\theta h)h,v\rangle
        \le\|F'(a+\theta h)\|\,|h|\,|v|.
\]
Делим на $|v|$ и полагаем $c=a+\theta h$.
\end{proof}
